% MCM/ICM Summary Sheet
% 摘要页，独立、自包含，不计入页数限制

% === 参赛信息（提交注意事项 1.1：第一页必须填写参赛组别与参赛编号） ===
\begin{center}
\begin{tabular}{|c|c|c|}
\hline
所属类别 & 2026年“华数杯”大学生数学建模竞赛 & 参赛编号 \\
\hline
本科生组 & & CM2601262 \\
\hline
\end{tabular}
\end{center}
\vspace{1em}

\begin{center}
{\Large\bfseries 面向算电协同的多目标调度优化研究}
\end{center}
\vspace{0.5em}
\begin{center}
{\bfseries 摘\quad 要}
\end{center}
\vspace{0.5em}

% === 第1段 背景与问题阐述（短句化：拆超长句，去"针对面向"拗口） ===
针对算电协同的多目标调度优化问题，我们以50\,000条AI工作负载轨迹与六区域电力市场数据为基础，围绕运行成本与碳排放控制，构建了从基础算力调度到多区域算-储-电协同的四层递进优化模型。模型覆盖需求预测、跨区迁移、储能协同与全要素统一求解四个子问题。

% === 第2段 数据处理说明与重要发现（第一人称"我们"，发现式预处理） ===
我们对数据执行检查、构造与探索三步预处理。检查发现，训练任务贡献约80\%的GPU需求，西部电价与碳强度则明显低于东部。这一差异为迁移降碳创造了空间。随后我们构造任务时间窗、时延可达集与逐时电价碳强度等特征，作为各调度模型的输入。

% === 第3段 问题一（加主语，拆主语漂移的4分句） ===
问题一中，我们先检验GPU需求的可预测性。自相关、非线性依赖、跨序列因果与机器学习四组检验一致表明，逐时需求在常用线性框架下不可预测，常数均值即近似最优预测。我们据此以区域利用率极差最小化为目标构建MILP模型。求解器在472.4秒内给出全局最优解。区域利用率极差下降25.2\%，任务更加分配均衡，538项任务全部按时完成。

% === 第4段 问题二（第一人称，迁移不牺牲QoS） ===
问题二中，我们以成本与碳排放为双目标，把延迟容忍任务迁往低价低碳区。我们通过容量感知贪心分配与滚动时域MILP组成的双层结构求解。76.4\%的任务发生跨区域迁移，碳排放下降32.7\%，成本仅上升2.0\%。迁移后的平均网络时延仅35.3 ms，远低于任务时延上限。A/B区保持空置是成本最优边界。强制启用15\%利用率，成本也只增加2.4\%。

% === 第5段 问题三（储能价值判断 + 碳排放刚性 + 波动冲突） ===
问题三中，我们引入储能与新能源，构建以运行成本最小化为目标的连续线性规划模型，同时量化两类收益来源。储能充放电贡献99.36~M元，购电时点优化贡献75.15~M元。两者合计降低174.51~M元，约7.3\%。储能自身的边际价值应以99.36~M元为准，否则会高估储能贡献。模型同时揭示出碳排放刚性。六区域碳排放下限介于$\varepsilon{=}0.957$与0.994之间，碳排放已逼近物理下限。峰谷套利还使净购电波动不降反升。

% === 第6段 问题四（统一解六指标 + 基荷预填精度量化） ===
问题四中，我们统一优化任务调度、储能与购售电。全耦合模型约有3660万候选变量，超出常规求解能力。我们通过借鉴核电不调峰的基荷政策，给出了新能源基荷预填策略，固定95.3\% GPU-hour的延迟容忍任务。变量规模从3660万压至223个，缩小约$10^5$倍。让我们惊喜的是，原本难以求解的问题现在约19秒即可完成求解。统一解成本2166.64~M元，较问题三再降2.1\%，碳排放1960.78~kt。按时完成率100\%，加权平均时延28.07 ms。配对对比实验把基荷预填的精度损失量化在0.88\%以内。大邻域搜索LNS还给出极端情形下次优间隙的保守估计约2.3\%。

% === 第7段 收束（第一人称，稳健性 + 整体自评 + 推广） ===
我们针对电价$\pm50\%$、新能源$\pm20\%$与碳约束档位分别开展敏感性分析，核心结论对参数波动稳健。整体来看，模型体系在降本、降碳、求解效率与服务质量之间取得了良好平衡，指标方向与量级均符合预期结果，我们对该方案的表现感到满意。整套框架以可量化精度损失换取可解性，可推广至多区域资源协同、碳约束与新能源消纳的调度场景。

\vspace{1em}
\noindent{\bfseries 关键词}\quad 算力调度；碳感知；储能协同优化；混合整数线性规划；新能源基荷预填